\scanpage{166}
\mspage{84}

\deleted{Notons que pour un espace homogène} Traduisons en termes de théorie des groupes, de façon tout à fait générale, la donnée suivante :

\[
\left\{\begin{aligned}
\text{a)}\quad&\text{espace homogène $S$ sous un groupe $G_0$, avec ``origine'' fixée $s_0\in S$}\\
\text{b)}\quad&\text{surgroupe $G\supset G_0$ opérant sur $S$ (en prolongeant l'action}\\
 &\text{précédente), avec $G_0$ distingué dans $G_0$ (soit $Z'=G/G_0$)}\\
\text{c)}\quad&\text{sousgroupe central $Z$ de $G$,}
\end{aligned}\right.\tag*{(26)}
\]

% slanted
\margin{\uline{NB} On ne préjuge pas, dans cette présentation, si les groupes $G_0,G$ opèrent fidèlement. On reviendra sur ce point par la suite\ldots}

en \unsure{essayant|exigeant}, dans la mesure du possible, de faire une description qui soit ``\unsure{proche|près|produite} du groupe $G_0$'', et de groupes qui se décrivent aisément en termes de $G_0$, \deleted{\unsure{de}} ses sous-groupes etc.

La donnée a) revient à la donnée du couple

\[ G_0\supset\widetilde U \tag*{(27)} \]

d'un groupe $G_0$, et d'un sous-groupe $\widetilde U$ $(=G_{0s_0})$, qui permet de \unsure{réconstruire|reconstruire} $(S,s_0)$ comme

\[ (S,s_0)\simeq(G_0/\widetilde U,\underset{\substack{\text{classe de }1_G\\\text{dans }G_0/\widetilde U}}{\dot{1}_G}). \tag*{(28)} \]

La donnée b) s'exprime, tout d'abord, comme celle d'un diagramme de groupes

\diagram{\begin{tikzcd}[ampersand replacement=\&,column sep=large]
(29)\qquad G_0\arrow[r,hook,no head, "\mathrm{dist.}"]\&G\&\\
\deleted{\unsure{G_0}}\widetilde U\arrow[u,hook,no head]\arrow[r,hook,no head]\&B'\arrow[u,hook,no head]\&(\mathrm{NB}\ B'=G_{s_0})
\end{tikzcd}}

``cartésien et cocartésien'', cartésien signifiant

\[ \widetilde U=G_0\cap B'\qquad(\text{donc }\widetilde U\text{ dist. dans }B') \tag*{(30)} \]

``cocartésien'' s'exprimant par une des deux conditions équivalentes

\[
\left\{\begin{aligned}
G_0/\widetilde U&\underset{\rm can}{\xrightarrow{\sim}}G/B' &\text{qui équivaut à}\\
B'/\widetilde U&\mathrel{\overset{\sim}{\text{---}}}G/G_0 &\text{(isom de groupes).}
\end{aligned}\right.\tag*{(31)}
\]

Le groupe $G$, extension d'un groupe

% archive sheet number: 165
